% Part A

% Q1. Implement the bottom-up strategy in Object Oriented Programming (OOP) development. What should be done first?

% (A) Define main()
% (B) Integrate complete system immediately
% (C) Develop and test small modules
% (D) Write documentation

% Q2. Determine the correct function call by applying default arguments:     void fun(int x, int y = 5);

% (A) fun(void)
% (B) fun()
% (C) fun(10)
% (D) fun(10,20,30)

% Q3. Demonstrate friend function concept to access private members. Select the correct declaration inside class.

% (A) friend void display();
% (B) static void display();
% (C) virtual display();
% (D) private display();

% Q4. Use pass by reference for object Student. Select the correct function prototype.

% (A) void fun(Student s)
% (B) void fun(Student &s)
% (C) void fun(Student)
% (D) void fun()

% Q5. Use pointer-to-pointer concept for integer variable x. Select the correct declaration.

% (A) int **p;
% (B) int *p;
% (C) int **x;
% (D) int *p;

% Q6. Utilize pointer to object for class Student. Identify the correct declaration.

% (A) Student s;
% (B) Student ptr;
% (C) Object Student;
% (D) ptr Student;

% Part B

% Q7.1 Implement the auto keyword and range-based for loop to display elements of an array and compute their total in C++.
%                                                       OR
% Q7.2 Develop a C++ program demonstrating function overloading to perform addition of two integers and two floating-point numbers.

% Q8.1 Apply encapsulation and data hiding concepts to design a Student class containing roll number and marks.
%      Implement member functions to set and display the values.
%                                                               OR
% Q8.2 Apply destructor concept in C++ by implementing a program that displays a message when an object is destroyed to illustrate object lifetime/scope.

% Q9.1 Apply dynamic memory management using new and delete operators by implementing a program that allocates memory for an integer and displays its value.
%                                                               OR
% Q9.2 Implement binary operator overloading using a member function to add two complex numbers represented as objects.

% Part C

% Q10.1 Apply range-based loops and string operations to develop a program that counts vowels, consonants, and digits in a given string.
%                                                           OR
% Q10.2 Construct a C++ program demonstrating interaction among multiple functions using different parameters passing techniques
%       (value, reference, pointer).

% Q11.1 Utilize friend class concept to implement a program where one class accesses private data of another class and explain the access mechanism.
%                                                                   OR
% Q11.2 Apply parameterized constructors and dynamic constructors to initialize objects containing dynamically allocated memory 
%       and demonstrate memory allocation through a program.

% Q12.1 Utilize operator overloading rules to implement unary operator (++) overloading for a class representing numeric values.
%       Write a program demonstrating the overloaded unary operator and explain the rules followed during operator overloading.
%                                                           OR
% Q12.2 Use relational operator overloading to compare objects of a user-defined class Distance. Develop the program and explain 
%       how relational operator overloading simplifies comparisons between objects.